\par
In regression the target column is a continuous variable. Meanwhile in
classification it consists of a target class (for example, "cat" or "dot", is it
"rainy" or "not rainy"), treated as binary outcomes. The outcome column
therefore changes its name to "label".

\par
Predictions are probabilistic, meaning that the model outputs the probability
for each "label". One can turn this probability into a decisive action or class,
through a decision threshold (could be $0.5$, but it varies based on the
necessity of the variable).

\par
In regression, we are looking for the line that is closest to all the points at
the same time. On the other hand in classification, we are looking for the most
reliable dividing line between the points, called the "decision boundary".

\par
To do this, a logistic regression is used instead (having a sigmoid shape),
being symmetric:

% TODO: include start

\begin{center}
	Logistic function: $\mathbb{R}$               $\rightarrow (0,1)$ \\
	probabilities($p$) $\leftarrow$ logistic($x$) $:= \frac{1}{1 + e^{-\kappa_1x +
					\kappa_0}}$
\end{center}

\par
Before, $\beta_i$ was found for the best linear fit, and in this case,
$\kappa_0$ and $\kappa_1$ are used to find the best logistic fit.
